% Source: https://artofproblemsolving.com/community/c3h1248201_tuff_stuff_geo_problem
% Stable identity: aops-topic-1248201-q01
Use coordinates with $O$ at the origin, $AB$ on the $x$-axis. The small circles have centres $(\pm16.5,0)$ and radius $16.5$. A circle with centre $(0,k)$ and radius $t$ that is internally tangent to the big circle and externally tangent to the small ones satisfies
\[|k|=33-t,\qquad\sqrt{16.5^2+k^2}=16.5+t.\]

\begin{asy}
usepackage("amsmath");
size(4.5cm);
defaultpen(fontsize(10pt));
draw(circle((0,0),33)); draw(circle((16.5,0),16.5)); draw(circle((-16.5,0),16.5));
draw(circle((0,22),11),blue); draw(circle((0,-22),11),blue);
draw((-33,0)--(33,0),gray); draw((0,-33)--(0,33),gray);
draw((0,22)--(16.5,0),red+dashed); dot((0,22)); dot((16.5,0));
label("$A$",(-33,0),W); label("$B$",(33,0),E); label("$C$",(0,33),N); label("$D$",(0,-33),S); label("$O$",(0,0),SW);
label("\mbox{}",(min(currentpicture,user=true)+max(currentpicture,user=true))/2,black); // keeps the text after the figure black when LaTeX typesets the labels inline
\end{asy}

Squaring the second and substituting $k^2=(33-t)^2$:
\[16.5^2+(33-t)^2=(16.5+t)^2\quad\Longrightarrow\quad1089-66t=33t\quad\Longrightarrow\quad t=\boxed{11}\text{ cm}.\]
(The two circles are centred at $(0,\pm22)$.)
